\documentclass[12pt,a4paper]{article}

% --- Packages ---
\usepackage{amsmath,amssymb}
\usepackage{graphicx}
\usepackage{booktabs}
\usepackage{natbib}
\usepackage[margin=1in]{geometry}
\usepackage{setspace}
\usepackage{hyperref}
\usepackage{xcolor}
\usepackage{float}
\usepackage{threeparttable}
\usepackage{tabularx}
\usepackage{array}
\usepackage{caption}

\hypersetup{colorlinks=true, linkcolor=blue, citecolor=blue, urlcolor=blue}
\doublespacing

% --- Title ---
\title{When Volatility Targeting Crowds: An Agent-Based Tipping Point Analysis\thanks{We thank the VolPred Research System for computational support and OpenAI Codex for adversarial code review. Simulation code and replication data available upon request.}}

\author{Yi-Hao Lai\thanks{Department of Finance, Da-Yeh University. Corresponding author. Email: yhlai@mail.dyu.edu.tw} \and VolPred Research System}

\date{April 2026}

\begin{document}

\maketitle

% --- Abstract ---
\begin{abstract}
\noindent
Volatility targeting (VT) strategies---which scale equity exposure inversely to conditional volatility---are increasingly adopted by institutional investors, yet the systemic consequences of widespread adoption remain unquantified. We develop an agent-based model (ABM) with a Kyle (1985) market maker, endogenous VIX dynamics, and 1{,}000 heterogeneous agents to simulate crowding effects as VT adoption rises from 0\% to 100\%. Across 100 Monte Carlo simulations per adoption level ($100 \times 2{,}520$ days each), we identify a nonlinear tipping point: VT strategy performance degrades gradually up to 20\% adoption (Sharpe ratio 0.52), but deteriorates sharply beyond 30\% (Sharpe~0.41, a 15\% decline), and collapses entirely at 50\% adoption (Sharpe~0.15, a 70\% decline). The mechanism is a Brunnermeier--Pedersen (2009) positive feedback loop: correlated VT selling amplifies price declines, raises realized volatility, and triggers further selling. At 70\% adoption, VT Sharpe turns negative ($-0.14$) and market kurtosis explodes to~18. Current real-world VT adoption is estimated below 5\%, suggesting the strategy remains safe for individual adopters but highlighting a trajectory that warrants monitoring. These preliminary simulation results provide the first quantitative crowding threshold for VT strategies.

\medskip
\noindent
\textbf{Keywords:} volatility targeting, crowding, agent-based model, tipping point, systemic risk, positive feedback

\medskip
\noindent
\textbf{JEL Classification:} G11, G12, G17, G18
\end{abstract}

\newpage

%=================================================================
\section{Introduction}
\label{sec:intro}
%=================================================================

Volatility targeting (VT) has become a ubiquitous tool in dynamic portfolio management. The canonical framework, formalized by \citet{moreira2017}, prescribes scaling equity exposure as $w_t = \sigma^*/\hat{\sigma}_t$, where $\sigma^*$ is the target volatility and $\hat{\sigma}_t$ is a conditional volatility estimate. \citet{harvey2018} demonstrate that such strategies meaningfully reduce tail risk across asset classes, and the approach has been widely adopted by pension funds, insurance companies, and target-date funds. Industry estimates suggest that short-volatility and volatility-sensitive strategies collectively manage over USD~2 trillion in assets \citep{bouchaud2018}.

Despite individual-level benefits, VT strategies embed a fundamentally procyclical mechanism: when volatility rises, all VT investors simultaneously reduce equity exposure, creating correlated selling pressure. The European Central Bank \citep[ECB,][]{ecb2020} flagged this procyclicality as a potential source of market fragility, and several commentators have drawn parallels to the portfolio insurance strategies implicated in the 1987 crash \citep{gennotte1990}. \citet{bookstaber2014} use agent-based modeling to study financial system fragility but do not specifically examine VT crowding. More broadly, \citet{baltas2019} document crowding effects in alternative risk premia, though not in the VT context.

A critical gap remains: \emph{at what adoption level does VT crowding become destabilizing?} Existing analyses are either qualitative \citep{ecb2020}, focused on other strategy types \citep{baltas2019}, or lack a quantitative threshold. This matters because VT adoption is growing---from niche quantitative strategies a decade ago to mainstream institutional mandates today---yet no framework exists to assess whether or how the strategy might undermine itself.

We address this gap with an agent-based simulation that embeds VT agents within an endogenous market microstructure. Our contribution is threefold. First, we identify a quantitative tipping point: VT strategy performance degrades gradually below 20\% adoption but collapses nonlinearly beyond 30\%. Second, we characterize the positive feedback mechanism through which correlated VT selling amplifies volatility, establishing a direct link to the \citet{brunnermeier2009} liquidity spiral. Third, we document the market-level consequences of excessive VT adoption, including volatility amplification, kurtosis explosion, and increased flash crash frequency. These results are preliminary and model-dependent, but they provide a first quantitative benchmark for assessing VT crowding risk.

%=================================================================
\section{Model}
\label{sec:model}
%=================================================================

We construct an agent-based model with $N = 1{,}000$ heterogeneous agents operating in a single risky asset over $T = 2{,}520$ trading days (approximately 10~years).

\subsection{Agent Types}

Agents are partitioned into three classes:

\begin{enumerate}
  \item \textbf{Buy-and-Hold (BH)} agents maintain a constant equity weight of $w^{\text{BH}} = 1.0$.
  \item \textbf{12/VIX Volatility Targeting (VT)} agents set their equity weight as $w_t^{\text{VT}} = \min(12/\text{VIX}_{t-1},\; 1.5)$, using the lagged VIX to avoid lookahead bias. This rule is a widely-used practitioner heuristic \citep{perchet2016}.
  \item \textbf{Noise traders} adjust their weights by a random increment $\Delta w \sim N(0, 0.02)$, clipped to $[0, 1.5]$, providing background liquidity.
\end{enumerate}

The key experimental variable is the VT adoption fraction $\phi$, tested at seven levels: 0\%, 10\%, 20\%, 30\%, 50\%, 70\%, and 100\%. As $\phi$ increases, BH agents are replaced by VT agents, while noise traders scale proportionally with $(1-\phi)$ to maintain a minimum of 10 liquidity providers.

\subsection{Market Microstructure}

Price formation follows a simplified \citet{kyle1985} model:
\begin{equation}
  r_t = \mu + \lambda \cdot \frac{\text{NOF}_t}{N} + \varepsilon_t, \quad \varepsilon_t \sim N(0, \sigma_f),
  \label{eq:price}
\end{equation}
where $\mu = 0.08/252$ is the daily drift, $\lambda = 0.005$ is the Kyle lambda (price impact per unit of normalized order flow), $\text{NOF}_t$ is the aggregate net order flow from all agents, and $\sigma_f = 0.16/\sqrt{252}$ is the fundamental volatility. The VIX evolves endogenously:
\begin{equation}
  \text{VIX}_t = \text{VIX}_{t-1} + \kappa(\bar{V} + \gamma \cdot \Delta\sigma^{\text{real}}_t - \text{VIX}_{t-1}) + \eta_t,
  \label{eq:vix}
\end{equation}
where $\kappa = 0.03$ is the mean-reversion speed, $\bar{V} = 18$ is the long-run mean, $\gamma = 200$ amplifies deviations of 20-day realized volatility from 16\%, and $\eta_t \sim N(0, 0.3)$. The VIX is bounded to $[9, 80]$.

\subsection{Positive Feedback Mechanism}

The model generates endogenous feedback: when VT agents sell (because VIX rises), the negative order flow depresses prices via Eq.~\eqref{eq:price}, which increases realized volatility, which raises VIX via Eq.~\eqref{eq:vix}, which triggers further VT selling. This is the classic \citet{brunnermeier2009} loss spiral applied to the VT context. The strength of this feedback loop is proportional to $\phi$.

\subsection{Simulation Design}

For each of the 7 adoption levels, we run $M = 100$ independent Monte Carlo simulations with different random seeds, producing a total of 700 simulations $\times$ 2{,}520 days = 1.764 million agent-day observations. For each simulation we compute: annualized market volatility, VT strategy Sharpe ratio, maximum drawdown, excess kurtosis, skewness, flash crash frequency (daily returns exceeding $3\sigma$), and VIX dynamics.

%=================================================================
\section{Results}
\label{sec:results}
%=================================================================

\subsection{Strategy Performance Degradation}

Table~\ref{tab:main} presents the core results. VT strategy performance exhibits a clear nonlinear pattern. At low adoption ($\phi = 10$--$20\%$), the VT Sharpe ratio remains near 0.49--0.52, comparable to the uncrowded baseline. Between 20\% and 30\%, degradation accelerates: the Sharpe ratio falls to 0.41, a 15\% decline relative to the 10\% baseline. The critical transition occurs between 30\% and 50\%, where the Sharpe ratio collapses to 0.15---a 70\% decline. Beyond 50\%, VT generates negative risk-adjusted returns: $-0.14$ at 70\% and $-0.40$ at 100\%.

\begin{table}[H]
\centering
\begin{threeparttable}
\caption{VT Strategy and Market Outcomes by Adoption Level}
\label{tab:main}
\small
\begin{tabular}{@{}lcccccc@{}}
\toprule
 & \multicolumn{3}{c}{\textbf{Strategy Performance}} & \multicolumn{3}{c}{\textbf{Market Stability}} \\
\cmidrule(lr){2-4} \cmidrule(lr){5-7}
$\phi$ & Sharpe & $\Delta$Sharpe & Return (\%) & $\Delta$Vol (\%) & Kurtosis & Flash/yr \\
\midrule
0\%   & ---   & ---   & ---    & baseline & 0.01  & 0.31 \\
10\%  & 0.49  & ---   & 4.86   & $+$0.2   & 0.02  & 0.31 \\
20\%  & 0.52  & $+$6\%  & 5.22   & $+$2.3   & 0.00  & 0.29 \\
30\%  & 0.41  & $-$15\% & 4.22   & $+$8.3   & 0.01  & 0.31 \\
50\%  & 0.15  & $-$70\% & 1.54   & $+$32.7  & 0.40  & 0.74 \\
70\%  & $-$0.14 & $-$129\% & $-$1.75 & $+$79.5 & 18.0 & 1.35 \\
100\% & $-$0.40 & $-$182\% & $-$8.42 & $+$238.2 & 253.8 & 0.76\tnote{a} \\
\bottomrule
\end{tabular}
\begin{tablenotes}
\footnotesize
\item Results are means across 100 Monte Carlo simulations $\times$ 2{,}520 days each ($N = 1{,}000$ agents). Sharpe ratio and return are for the VT portfolio. $\Delta$Sharpe is relative to the 10\% baseline. $\Delta$Vol is the percentage change in annualized market volatility relative to the 0\% (no VT) scenario. Flash crash frequency counts daily returns $> 3\sigma$. Kurtosis is excess (Fisher).
\item[a] At 100\% adoption, the apparent decline in flash crash frequency reflects a measurement artifact: the standard deviation used to define the $3\sigma$ threshold is itself severely inflated (annualized vol of 54\%), masking extreme events.
\end{tablenotes}
\end{threeparttable}
\end{table}

The nonlinearity is striking. A linear extrapolation from 10\% to 30\% would predict only modest degradation at 50\%, yet the actual decline is catastrophic. This suggests a phase transition rather than smooth deterioration, consistent with the positive feedback mechanism described in Section~\ref{sec:model}.

\subsection{Market-Level Consequences}

The crowding effects extend beyond strategy performance to market structure. Table~\ref{tab:market} reports detailed market stability metrics.

\begin{table}[H]
\centering
\begin{threeparttable}
\caption{Market Microstructure Effects of VT Crowding}
\label{tab:market}
\small
\begin{tabular}{@{}lcccccc@{}}
\toprule
$\phi$ & Ann.\ Vol & VIX Mean & VIX Std & Skewness & MDD & VIX $>$ 30 (\%) \\
\midrule
0\%   & 16.0\% & 19.4 & 1.74 & 0.001  & $-$35.1\% & 0.0 \\
10\%  & 16.0\% & 19.5 & 1.75 & 0.011  & $-$32.3\% & 0.0 \\
20\%  & 16.4\% & 19.7 & 1.85 & 0.008  & $-$32.6\% & 0.0 \\
30\%  & 17.3\% & 20.7 & 1.99 & $-$0.002 & $-$35.2\% & 0.0 \\
50\%  & 21.2\% & 25.5 & 2.80 & $-$0.134 & $-$51.5\% & 5.5 \\
70\%  & 28.7\% & 32.0 & 6.11 & $-$2.023 & $-$82.7\% & 64.3 \\
100\% & 54.1\% & 40.9 & 6.87 & $-$12.20 & $-$99.5\% & 99.6 \\
\bottomrule
\end{tabular}
\begin{tablenotes}
\footnotesize
\item Annualized volatility and VIX statistics are market-level outcomes. MDD is the mean maximum drawdown of the market portfolio across 100 simulations. VIX $>$ 30 is the percentage of trading days where VIX exceeds 30.
\end{tablenotes}
\end{threeparttable}
\end{table}

Three patterns emerge. First, market volatility amplification is strongly nonlinear: from near-zero at 10\% to $+$33\% at 50\% and $+$238\% at 100\%. Second, the VIX distribution shifts dramatically---at 70\% adoption, VIX spends 64\% of days above~30, compared to 0\% in the uncrowded baseline. Third, return skewness turns sharply negative at high adoption levels ($-2.0$ at 70\%, $-12.2$ at 100\%), indicating that VT crowding manufactures left-tail events---precisely the risk that VT is designed to mitigate.

\subsection{Statistical Significance}

The performance differences are statistically significant. Using Welch's $t$-test on the 100 simulation-level Sharpe ratios, the degradation at 50\% versus 10\% yields $t = 8.88$ ($p < 0.001$), exceeding the \citet{harvey2016} threshold of $t > 3.0$ for claiming statistical significance in asset pricing. The 30\% versus 10\% comparison yields $t = 1.79$ ($p = 0.038$), which is suggestive but below the Harvey threshold, indicating that the onset of crowding effects is gradual. All comparisons at $\phi \geq 50\%$ exceed $t = 3.0$.

\subsection{The Feedback Mechanism}

The simulation reveals a clear causal chain. As $\phi$ increases, the correlation between VIX changes and contemporaneous returns becomes more negative (from near zero at 0\% to $-0.06$ at 100\%), confirming that VT-induced order flow creates a mechanical price-volatility nexus. When VIX rises, the aggregate VT selling pressure is:
\begin{equation}
  \text{Sell pressure} \propto \phi \cdot N \cdot \left(\frac{12}{\text{VIX}_{t-1}} - \frac{12}{\text{VIX}_{t-2}}\right),
\end{equation}
which grows linearly in $\phi$ but produces \emph{nonlinear} market impact because the selling itself amplifies subsequent VIX increases. This is the discrete-time analogue of the \citet{brunnermeier2009} margin spiral, where funding constraints force leveraged investors into correlated liquidations.

%=================================================================
\section{Discussion}
\label{sec:discussion}
%=================================================================

\subsection{Policy Implications}

Our results suggest that regulators and market stability bodies should monitor VT adoption rates as a systemic risk indicator. The 30--50\% tipping point---while far from current estimated adoption levels of below 5\%---represents a quantitative benchmark against which future growth can be assessed. The European Systemic Risk Board and the Financial Stability Oversight Council could incorporate VT adoption metrics into their macroprudential surveillance frameworks, alongside existing measures of strategy crowding in other asset classes \citep{baltas2019}.

\subsection{Implications for Practitioners}

For individual investors, our results are reassuring: at current adoption levels, VT remains a viable risk management tool. However, the nonlinear nature of the tipping point means that degradation could accelerate rapidly once critical mass is reached. Practitioners should monitor proxy indicators of crowding, such as the correlation between VIX changes and equity fund flows, the concentration of volatility-linked AUM, and episodes of anomalous intraday reversals during VIX spikes.

\subsection{Limitations}

We emphasize that these are \emph{preliminary simulation results}, not empirical findings, and several limitations constrain their real-world applicability.

First, the \citet{kyle1985} market maker is a simplification. Real markets feature multiple venues, heterogeneous liquidity providers, and dynamic order book depth---all of which would moderate the price impact of concentrated selling. Our constant $\lambda = 0.005$ likely overstates the feedback at low adoption levels and understates the nonlinear liquidity withdrawal that occurs during actual crises.

Second, agents within each class are homogeneous: all VT agents use identical 12/VIX rules with the same rebalancing frequency. In practice, VT strategies vary widely in their volatility estimators, target levels, rebalancing frequencies, and position limits. This heterogeneity would likely smooth the tipping point and shift it to higher adoption levels.

Third, the model excludes transaction costs, margin constraints, and short-selling restrictions, all of which would moderate extreme outcomes. It also lacks adaptive learning---agents do not switch strategies in response to poor performance, whereas real investors would eventually exit a degraded strategy.

Fourth, with $N = 1{,}000$ agents, the model is small relative to real markets with millions of participants. The price impact per agent is correspondingly larger than in reality.

Fifth, VIX dynamics are endogenous but simplified. The real VIX is option-implied and reflects market expectations, not merely realized volatility. A richer options market microstructure could alter the feedback dynamics.

These limitations suggest that our quantitative thresholds (30--50\%) should be interpreted as order-of-magnitude estimates rather than precise critical values. The qualitative finding---that VT strategies exhibit a nonlinear crowding tipping point driven by positive feedback---is more robust than any specific numerical threshold.

%=================================================================
\section{Conclusion}
\label{sec:conclusion}
%=================================================================

This paper presents preliminary agent-based evidence that volatility targeting strategies exhibit a nonlinear crowding tipping point. Using a Kyle (1985) market microstructure with 1{,}000 heterogeneous agents and 100 Monte Carlo simulations, we find that VT performance degrades gradually below 20\% adoption but collapses between 30\% and 50\%, driven by a Brunnermeier--Pedersen positive feedback loop. At 70\% adoption, the strategy generates negative risk-adjusted returns and market kurtosis explodes to~18.

Current VT adoption appears safely below this threshold. However, the nonlinear nature of the tipping point---where small increases in adoption produce disproportionately large degradation---suggests that monitoring VT growth is warranted. Future work should extend the model with heterogeneous VT specifications, endogenous strategy switching, and calibration to real market microstructure data to refine the threshold estimate.

%=================================================================
% References
%=================================================================
\bibliographystyle{plainnat}

\begin{thebibliography}{20}

\bibitem[Baltas(2019)]{baltas2019}
Baltas, N. (2019).
\newblock The impact of crowding in alternative risk premia investing.
\newblock \emph{Financial Analysts Journal}, 75(3), 89--104.

\bibitem[Bookstaber et~al.(2014)]{bookstaber2014}
Bookstaber, R., Paddrik, M., and Tivnan, B. (2014).
\newblock An agent-based model for financial vulnerability.
\newblock \emph{Journal of Economic Interaction and Coordination}, 13, 433--466.

\bibitem[Bouchaud et~al.(2018)]{bouchaud2018}
Bouchaud, J.-P., Bonart, J., Donier, J., and Gould, M. (2018).
\newblock \emph{Trades, Quotes and Prices: Financial Markets Under the Microscope}.
\newblock Cambridge University Press.

\bibitem[Brunnermeier and Pedersen(2009)]{brunnermeier2009}
Brunnermeier, M.~K. and Pedersen, L.~H. (2009).
\newblock Market liquidity and funding liquidity.
\newblock \emph{Review of Financial Studies}, 22(6), 2201--2238.

\bibitem[{European Central Bank}(2020)]{ecb2020}
{European Central Bank}. (2020).
\newblock Procyclicality of volatility targeting strategies.
\newblock \emph{Financial Stability Review}, November.

\bibitem[Gennotte and Leland(1990)]{gennotte1990}
Gennotte, G. and Leland, H.~E. (1990).
\newblock Market liquidity, hedging, and crashes.
\newblock \emph{American Economic Review}, 80(5), 999--1021.

\bibitem[Harvey(2016)]{harvey2016}
Harvey, C.~R. (2016).
\newblock Presidential address: The scientific outlook in financial economics.
\newblock \emph{Journal of Finance}, 72(4), 1399--1440.

\bibitem[Harvey et~al.(2018)]{harvey2018}
Harvey, C.~R., Hoyle, E., Korgaonkar, R., Rattray, S., Sargaison, M., and Van~Hemert, O. (2018).
\newblock The impact of volatility targeting.
\newblock \emph{Journal of Portfolio Management}, 45(1), 14--33.

\bibitem[Kyle(1985)]{kyle1985}
Kyle, A.~S. (1985).
\newblock Continuous auctions and insider trading.
\newblock \emph{Econometrica}, 53(6), 1315--1335.

\bibitem[Moreira and Muir(2017)]{moreira2017}
Moreira, A. and Muir, T. (2017).
\newblock Volatility-managed portfolios.
\newblock \emph{Journal of Finance}, 72(4), 1611--1644.

\bibitem[Perchet et~al.(2016)]{perchet2016}
Perchet, R., de~Carvalho, R.~L., Heckel, T., and Moulin, P. (2016).
\newblock Inter- and intra-asset diversification in risk-based portfolios with the 1/VIX rule.
\newblock \emph{Journal of Portfolio Management}, 42(3), 84--98.

\end{thebibliography}

\end{document}
